\section{Alimentación cruzada del control ambiental ordinario: ENV-40}
\label{sec:ambiente-env40}
\textbf{Selección y alcance.} Se seleccionan dos Phoenix Contact \texttt{QUINT-ORING/24DC/2X10/1X20}, artículo \texttt{2320173}, uno para el riel ambiental de cada banco. Cada módulo recibe un positivo de la fuente A y otro de la fuente B \texttt{QUINT4-PS/1AC/24DC/10}, artículo \texttt{2904601}, ya seleccionadas; conserva una salida separada por banco. No se añaden dos fuentes ni se reemplazan las QUINT de los puestos. El módulo publica entrada 18--28 V DC, dos entradas nominales de 10 A, salida de 10 A en redundancia o 20 A en aumento de potencia, caída nominal 0,1 V y pérdida máxima publicada de 2 W a salida de 20 A. La capacidad de redundancia aplicable es 10 A; no se adopta 20 A como capacidad mantenida ante una entrada retirada. ACB distribuye corriente entre entradas; no convierte el conjunto en fuente ininterrumpible ni mantiene cargas AC \parencite[pp.~1--3]{lote40-oring-ficha}\parencite{lote40-oring-catalogo}.

Se trasladan al riel ambiental respaldado de cada banco sólo CPU y módulos WAGO completos, dos DMT132 y sus lazos/DI (25,9776 W de reserva ENV-34), dos RVAZ4-24A (12 W de reserva ENV-31) y dos DPT de caudal (2,4 W): 40,3776 W por banco, 80,7552 W entre ambos. La alimentación de campo AI/AO y las referencias locales de esos receptores se trasladan conjuntamente; no se permite una AO energizada desde el riel respaldado con actuador o transmisor cuyo suministro quede en el camino retirado. El consumo de bus del 750-479 permanece incluido dentro de los 24 W por nodo; las fuentes TEN de diagnóstico y sus receptores externos no migran por esa inclusión. Los 72,8 W de permisos/SEL, 195,2 W de TEN/PTCB y 4 W de interfaces conservan reparto de 136 W por fuente, sin afirmar que el diagnóstico del banco retirado siga disponible. La reserva independiente de H$_2$ de ECI-29 (27,84 W continuos y 67,2 W de arranque) no pertenece a esta base ni se incorpora al OR ambiental. La seguridad de gas no depende de esta CPU.

\textbf{Cálculo de continuidad de suministro.} Sobre la base reconciliada de 352,7552 W DC se reservan provisionalmente 2 W por OR, 4 W adicionales. Es reserva propia fundada en la pérdida nominal publicada, pendiente de confirmar consumo y pérdidas a carga parcial/arranque; no es un nuevo máximo OEM del conjunto. Con reparto simétrico, cada fuente sostiene 136+40,3776+2=178,3776 W. Al retirar una fuente, la superviviente sostiene su carga no migrada y ambos rieles ambientales: $136+80,7552+4=220,7552$ W, o 9,198133 A a 24 V. Quedan 19,2448 W respecto a 240 W nominales; cada entrada OR que permanece soporta aproximadamente $(40,3776+2)/24=1,765733$ A. No se conserva el margen anterior de 63,6224 W como margen N+1, ni se migra la totalidad de 356,7552 W a una fuente de 240 W. Aplicar un margen adicional del 20\,\% exigiría 264,90624 W y no cabe en la selección actual. Las cotas CPU, arranques de actuadores/accesorios y cargas adicionales requieren cotejo antes de liberar el suministro.

La fuente publica 240 W/10 A nominales, reducción de 2,5\,\%/K entre 60 y 70\,$^{\circ}$C y reducción de entrada por debajo de 100 V AC. A 65\,$^{\circ}$C su capacidad nominal queda en 210 W, insuficiente para el estado anterior. Se conserva el armario interior de proyecto a 24\,$^{\circ}$C, seco y sin condensación, y se establece límite de diseño de 60\,$^{\circ}$C para las fuentes y los OR; su disipación/ventilación y vigilancia deben demostrarlo. Se alimentan a 230 V desde los caminos protegidos vigentes. No se usa Static Boost, Dynamic Boost o SFB para completar capacidad continua, ni eficiencia típica de plena carga para calcular consumos a esta carga parcial. Los 4 W nuevos se incorporan a conversión/pérdidas de fuente, autonomía y disipación; no se suman también como 4 W AC \parencite{lote40-quint-catalogo}.

\textbf{Referencia DC, tensión y protecciones.} Esta conexión sustituye, para los receptores ambientales enumerados, la alimentación exclusiva de su camino establecida en ENV-34. Los positivos A/B permanecen desacoplados dentro de cada OR; no se puentean sus entradas. El OR no proporciona aislamiento galvánico entre ambos suministros. Sus retornos requieren una referencia 0 V común intencional en el punto de distribución definido, con recorrido de retorno propio por cada alimentación y unión de protección/PE documentada. Los circuitos de señal de cada banco conservan su cableado propio; los retornos de campo ambientales comparten la referencia secundaria común seleccionada. No se atribuye aislamiento entre esas referencias 0 V. Las referencias de medida aisladas de TEN/LEM permanecen independientes entre bancos; pantallas, AO o masas no sustituyen un retorno. La supervivencia de un retorno roto, cortocircuito de salida, fallo interno de OR o CPU única no está acreditada por desacoplar positivos. La protección de entrada por fuente y salida por banco debe seleccionar dispositivo, corriente, cable y selectividad antes de habilitar la transferencia; los PTCB de 1 A destinados a permisos no se reutilizan como protección de toda la carga ambiental.

Se fija ajuste de fuente a 24 V. Como cálculo de tendido propio se reserva cobre de sección efectiva mínima 1,5 mm$^2$, hasta 10 m de ida por tramo fuente--OR y OR--receptores, resistividad 0,0175 $\Omega$ mm$^2$/m. Cada ida/retorno de 20 m resulta en 0,233333 $\Omega$; a 1,765733 A cae 0,412004 V, y dos tramos agregan 0,824009 V. Es presupuesto resistivo a las condiciones indicadas, pendiente de temperatura, conexiones, cable comercial y verificación de ampacidad/protección. La caída nominal OR de 0,1 V deja 23,075991 V desde 24 V ideales antes de desviación de fuente y transitorios, por encima de 20,4 V mínimos Regin y 20 V exigidos para conservar el presupuesto de lazo DMT de ENV-34. No se presenta esa resta como tensión mínima garantizada durante retirada/arranque. La protección OEM de sobretensión de fuente/OR hasta 32 V no garantiza los máximos de 27,6 V Regin y 28 V DMT/OR; falta protección compatible y presupuesto de tensión máximo. Las pruebas de tensión se realizan en bornes del receptor y conservan las cotas de carga/entrada AO aún pendientes.

\textbf{Implantación, estados y recepción.} Se instalan dos módulos de 32 por 130 por 125 mm en riel DIN horizontal, fuera del retorno con gas. Se reserva por módulo 15 mm hacia componentes activos y 50 mm arriba/abajo, sin reducir mantenimiento; el UTA107/30 de montaje está incluido en el módulo y no constituye una tercera partida. Se contemplan cuatro circuitos de entrada y dos de salida, con positivos y retornos identificados. Los contactos 13--14 Redundancy OK y 23--24 ACB OK permiten supervisión ordinaria; esta variante no detecta cable roto a la salida. Los dieciséis DI WAGO de ENV-34 ya están asignados: no se conectan cuatro señales nuevas a canales supuestamente libres ni se declara esa supervisión instalada antes de seleccionar una ampliación/interfaz. La caída de una entrada genera operación degradada del suministro, sin autorizar por sí misma recarga. El dato LEM de un banco sin su fuente externa se invalida aunque la CPU continúe encendida. Fallo de CPU, OR, referencia o tensión del receptor enclava condición de control no disponible y deja la decisión de retirada de recarga/calor y ventilación segura a la cadena independiente. El retorno de energía no impone rearme automático.

La recepción prevista registra retirada de A y de B, tensión/corriente en cada receptor, arranque conjunto, alarmas, pérdida de retornos, corto de ramal y recuperación; comprueba ausencia de reinicio, salida analógica transitoria indebida y datos retenidos como válidos. Se contrasta con ventiladores AC, bomba, Carrier y transferencia hidráulica de ENV-37: sostener mando DC no demuestra continuidad de esos motores, frío o caudal. Permanecen abiertos protección/selectividad/OVP, máximo y arranque de cargas, tierra/retornos, alarmas, fuente/serpentín/primario y seguridad independiente; RT-06.03, RT-06.09 y RT-06.13 no se declaran completamente desarrollados por esta selección parcial de continuidad.
